\newpage
\section{Discussion}

\subsection{A Trust Hierarchy for Multi‑Prover Verification}

The benchmark and the follow‑up verification underline that not all provers
are equally trustworthy for every task.  We formalise this observation by
assigning each prover a \emph{trust level} $\tau \in [0,1]$ for a given
framework $f$ and observable type $T$:

\[
\tau(\text{prover}, f, T) \in \{0, \tfrac{1}{2}, 1\},
\]

where $1$ denotes ``sound and fully automatic'', $\tfrac{1}{2}$ denotes
``sound under specific conditions (e.g., only for integers)'', and $0$
denotes ``not sound / placeholder''.  The certificate's \texttt{verified\_by}
list is then filtered to include only provers with $\tau \ge \tfrac{1}{2}$
and, for $\tau = \tfrac{1}{2}$, a metadata field records the restriction.
Currently:

\begin{itemize}
    \item $\tau(\text{Z3}, \text{boolean}, \text{int}) = 1$,
          $\tau(\text{Z3}, \text{boolean}, \text{set}) = 1$;
    \item $\tau(\text{Lean}, \text{boolean}, \text{int}) = \tfrac{1}{2}$,
          $\tau(\text{Coq}, \text{boolean}, \text{int}) = \tfrac{1}{2}$;
    \item all other combinations have $\tau = 0$.
\end{itemize}

The assignment of trust levels is not based solely on the benchmark
results reported in Section~4.  Z3 is an industrial‑strength SMT
solver whose soundness for quantifier‑free arithmetic has been
independently established~\cite{moura2008z3}; SymPy's propositional
simplification is, by construction, unable to detect existential
counterexamples.  Lean and Coq are interactive provers with small,
trusted kernels, but their automation for our specific domain is
still under development.  Hence the trust hierarchy reflects both
empirical evidence and the fundamental design of each tool.

This hierarchy is not static; as the Lean and Coq backends mature (e.g., by
implementing a cardinality check for sets), their trust levels can be raised
and automatically reflected in future certificates without changing the
client code.

\subsection{Comparative Analysis of Theorem Provers}

% figures/prover_comparison_table.tex
\begin{table}[htbp]
\centering
\caption{Comparative strengths and weaknesses of the integrated theorem provers.}
\label{tab:prover_comparison}
\begin{tabular}{p{2.2cm} p{3.5cm} p{3.5cm}}
\toprule
\textbf{Prover} & \textbf{Strengths} & \textbf{Weaknesses} \\
\midrule
SymPy & Pure Python, fast logical simplification, no external dependencies. & Limited to propositional logic; cannot handle quantifiers or complex arithmetic. \\
Z3 & Powerful SMT solving (arithmetic, bit-vectors, quantifiers); high automation. & External solver required; performance degrades on large theories; proofs not human-readable. \\
Lean 4 & Dependent type theory; highest assurance; large mathlib. & Requires manual proof effort; automation limited; steep learning curve. \\
Coq & Mature ecosystem; strong automation (Ltac); widely used in formal verification. & Steep learning curve; automation can be brittle; requires external compilation. \\
\bottomrule
\end{tabular}
\end{table}

For the specific task of atomicity verification, the trade‑offs are clear:
SymPy is fast but unsound; Z3 is the workhorse for automated, sound proofs;
Lean and Coq offer the highest potential assurance but require substantially
more engineering to reach full automation.  This division of labor mirrors
the design philosophy of the \Apeiron{} Framework itself: lightweight
heuristics guide the system, while formal proofs are reserved for safety‑
critical claims.

\subsection{Translation Semantics and the Semantic Gap}

A fundamental challenge in multi‑prover systems is ensuring that the
property being proved is preserved across translation.  Let $P$ be a
target property of an observable $o$, and let $\llbracket P \rrbracket_L$
be its encoding in prover $L$'s logic.  We say that the translation is
\emph{sound} if
\[
\forall o,\; \llbracket P \rrbracket_L \text{ is provable in } L
\implies P(o) \text{ holds}.
\]

Currently, soundness is manually validated for each prover–framework–
type combination (Section~4.2).  A formally verified translation layer,
built around an abstract syntax tree with certified serialisers, would
elevate the system to the highest assurance level and is a central item on
our research agenda.

A stronger approach—planned for the next iteration—is to define the
atomicity criteria \emph{directly} in the language of the interactive
provers.  For example, the Boolean atomicity of integers would be
expressed as a Lean predicate
\begin{lstlisting}[language=Python]
def isAtom (n : Nat) : Prop := n = 1
\end{lstlisting}
together with a proof that this definition is equivalent to the
divisibility‑based one.  The Python layer would then only supply the
concrete value \texttt{n} to be substituted.  This eliminates the
trusted code base in the generator and makes the semantic preservation
a theorem inside the prover rather than an implicit assumption.

\subsection{Toward an AST‑Based Translation Layer}

To close the semantic gap, we are developing a translation layer
centred on an \emph{abstract syntax tree} (AST).  In this design,
the \texttt{AtomicityTheoremGenerator} would no longer perform string
interpolation; instead, it would build a language‑agnostic AST
representing the atomicity claim.  Each prover backend would then
be a \emph{verified printer} that walks the AST and emits the
corresponding prover‑native code.  Because the AST nodes carry
precise semantic information (e.g., \texttt{Forall} for universal
quantification, \texttt{BinaryOp} for divisibility), a machine‑
checkable proof can be carried out \emph{once} for each printer,
ensuring that the generated code faithfully captures the original
claim.  This approach is inspired by the CompCert verified compiler
and by the DeepSpec consortium's methodology~\cite{deepspec2016,leroy2009coq}.
We estimate that the Boolean backend for integers can be migrated
to an AST‑based printer within a few person‑months; the other
frameworks will follow as their Lean/Coq formalisations mature.

\subsection{Scalability and Compositional Verification}

The scalability gap between verifying a single observable and verifying an
entire knowledge graph remains the primary obstacle.  We are investigating
compositional proof rules that allow the atomicity of a complex structure
(e.g., a functional cluster in Layer~3) to be derived from the atomicity
of its constituents, potentially reducing the proof burden from exponential
to linear in the number of components.

\subsection{Implications for AI Alignment}

The ability to produce machine‑checkable certificates of foundational
properties directly addresses a core alignment concern: if an AI builds its
own ontology, we need guarantees that its ``atoms'' are not arbitrary or
contradictory.  The JSON certificates, together with the trust hierarchy,
enable third‑party auditors to verify the system's internal claims without
trusting the AI itself—a crucial property for regulatory acceptance.

\subsection{Limitations and Future Work}

\begin{itemize}
    \item \textbf{Lean/Coq coverage}: Only Boolean atomicity for integers
          is currently trustworthy.  Extending the backends to sets,
          strings, and the other three frameworks is ongoing.
    \item \textbf{Kolmogorov proxy}: The zlib‑based information‑theoretic
          test is a heuristic; a true formal proof of incompressibility is
          impossible in general.
    \item \textbf{Certificate revocation}: When a prover's trust level is
          lowered (e.g., a bug is discovered), existing certificates should
          be re‑evaluated.  We plan to implement a versioned certificate
          store with automatic re‑verification.
\end{itemize}